\section*{अध्याय \ref{ch:b1:layered-earth} --- स्तरित पृथ्वी}

\begin{solution}{exo:b1:layered-earth:1}
आयतन-अंश त्रिज्या-अनुपातों के घन हैं: कोर के लिए $(3480/6371)^3 = 0.1630$, अर्थात् \qty{16.3}{\percent}; \omterm{def:b1:layered-earth:layers}{आंतरिक कोर} के लिए
$(1221.5/6371)^3 = 0.00705$, \qty{0.70}{\percent}।
\end{solution}

\begin{solution}{exo:b1:layered-earth:2}
ऊपर: $\Phi = 13.7166^2 - \tfrac43\times7.26466^2 = 188.145 - 70.367 =
\qty{117.78}{km^2.s^{-2}}$, $K = \rho\Phi = \qty{5566.45}{kg.m^{-3}}\times
\qty{117.78e6}{m^2.s^{-2}} = \qty{655.6}{GPa}$। नीचे: $\Phi = 8.06482^2 =
\qty{65.04}{km^2.s^{-2}}$, $K = 9903.49\times\num{65.04e6} = \qty{644.1}{GPa}$।
बल्क मापांक लगभग नहीं बदलता (\qty{-1.7}{\percent})। जो बदलता है वह अपरूपण-मापांक है, $\mu = \rho v_S^2 = \qty{293.8}{GPa}$
ऊपर और नीचे $0$ (एक द्रव), और घनत्व, जो \qty{78}{\percent} बढ़ता है: इसीलिए $v_P$ का गिरना $K$ के लगभग समान रहने के बावजूद होता है।
\end{solution}

\begin{solution}{exo:b1:layered-earth:3}
$\sin i_c = 6.2/8.1 = 0.7654$, $i_c = \ang{49.95}$। क्रांतिक दूरी
$2\times30\times\tan\ang{49.95} = 60\times1.1906 = \qty{71.4}{km}$ है। \omterm{def:b1:layered-earth:head-wave}{क्रॉसओवर
दूरी} $60\sqrt{14.3/1.9} = 60\times2.743 = \qty{164.6}{km}$ है।
\end{solution}

\begin{solution}{exo:b1:layered-earth:4}
$g = Gm/r^2 = \num{6.674e-11}\times\num{1.94e24}/(\num{3.48e6})^2 =
\qty{10.69}{m.s^{-2}}$, पृष्ठ-गुरुत्व से \qty{9}{\percent} अधिक। कोर की ओर नीचे जाने पर त्रिज्या का वर्ग अपने पृष्ठ-मान का
$(3480/6371)^2 = 0.298$ रह जाता है, जबकि आकर्षित करने वाला द्रव्यमान कुल का केवल $1.94/5.972 = 0.325$ रह जाता है,
क्योंकि कोर द्रव्यमान का एक-तिहाई रखता है: $0.325/0.298 = 1.09$।
\end{solution}

\begin{solution}{exo:b1:layered-earth:5}
स्रोत के निकट $v_1 = 60/10.0 = \qty{6.0}{km.s^{-1}}$। दूर
$v_2 = (300 - 200)/(45.8 - 33.3) = 100/12.5 = \qty{8.0}{km.s^{-1}}$; अंतःखंड-समय
$t_i = 33.3 - 200/8.0 = \qty{8.3}{s}$ है। तब $\sin i_c =
0.75$, $\cos i_c = 0.6614$, और $h = t_i v_1/(2\cos i_c) = 8.3\times6.0/1.3229
= \qty{37.6}{km}$। जाँच: $x_c = 2\times37.6\times\sqrt{14/2} = \qty{199}{km}$,
जहाँ दोनों तरंगें $199/6.0 = 199/8.0 + 8.3 = \qty{33.2}{s}$ पर पहुँचती हैं।
\end{solution}

\begin{solution}{exo:b1:layered-earth:6}
एकसमान $\rho$ के लिए, $I = \tfrac{8\pi}{3}\rho\int_0^R r^4\,\mathrm dr =
\tfrac{8\pi}{15}\rho R^5$ और $M = \tfrac43\pi\rho R^3$ है, अतः
$I/MR^2 = (8/15)/(4/3) = 0.4$। दो परतें ($r_c = \qty{3.48e6}{m}$,
$R = \qty{6.371e6}{m}$): $M = \tfrac43\pi\bigl(11000\,r_c^3 + 4450\,(R^3 -
r_c^3)\bigr) = \qty{5.977e24}{kg}$, मापे गए द्रव्यमान से \qty{0.08}{\percent} अधिक; $I = \tfrac{8\pi}{15}\bigl(11000\,r_c^5 + 4450\,(R^5 - r_c^5)\bigr)$
से $I/MR^2 = 0.3457$ मिलता है। द्रव्यमान सही है, पर गुणांक \qty{4.5}{\percent} अधिक है: वास्तविक \omterm{def:b1:layered-earth:layers}{प्रावार} एकसमान नहीं है,
उसका घनत्व गहराई के साथ बढ़ता है, जो द्रव्यमान को भीतर की ओर सरकाता है।
\end{solution}

\begin{solution}{exo:b1:layered-earth:7}
$\Phi = 12.16^2 - \tfrac43\times6.673^2 = 147.87 - 59.37 =
\qty{88.49}{km^2.s^{-2}}$। गहराई $z$ के साथ,
$\mathrm d\rho/\mathrm dz = \rho g/\Phi = 4860\times9.9/\num{88.49e6} =
\qty{5.44e-4}{kg.m^{-4}}$, अर्थात् प्रति किलोमीटर \qty{0.544}{kg.m^{-3}}, या प्रति \qty{1000}{km} \qty{0.54}{g.cm^{-3}}। PREM इन
\qty{1000}{km} पर $5.357 - 4.844 =
0.513$ बढ़ता है, \qty{6}{\percent} कम: ऊपर परिकलित प्रवणता पूरे अंतराल पर लागू की गई है, जबकि $\Phi$ गहराई के साथ
बढ़ता है और प्रवणता घटती है।
\end{solution}

\begin{solution}{exo:b1:layered-earth:8}
$P \approx \bar\rho g z = 3570\times9.9\times\num{670e3} = \qty{23.7}{GPa}$
(PREM: \qty{23.9}{GPa})। यह वह दाब है, लगभग \qty{24}{GPa}, जिस पर रिंगवुडाइट टूटकर ब्रिजमैनाइट और फ़ेरोपेरिक्लेज़ बनाता है:
\qty{660}{km} असांतत्य।
\end{solution}

\begin{solution}{exo:b1:layered-earth:9}
(क)~द्रव लोहा: $10.90\times1.05 = 11.4$ से $10.90\times1.10 =
\qty{12.0}{g.cm^{-3}}$। (ख)~$12.764/12.166 = 1.049$: PREM में \qty{4.9}{\percent} की छलांग, हाल के अध्ययनों के
\qty{4.5}{\percent} से थोड़ी अधिक। (ग)~अकेला जमना ठोस को द्रव से केवल थोड़ा घना बनाता; बड़ी छलांग का अर्थ है कि ठोस
में द्रव की तुलना में कम हल्के तत्व हैं: वे द्रव को पसंद करते हैं, और \omterm{def:b1:layered-earth:discontinuity}{आंतरिक-कोर सीमा} पर जमना उनमें से अधिकांश को पीछे छोड़ देता है।
\end{solution}

\begin{solution}{exo:b1:layered-earth:10}
माना $x_* = 2h\tan i_c$। तब $\sqrt{x_*^2 + 4h^2} = 2h/\cos i_c$ और
$t_r(x_*) = 2h/(v_1\cos i_c)$। \omterm{def:b1:layered-earth:head-wave}{हेड तरंग} देती है
$t(x_*) = 2h\tan i_c/v_2 + 2h\cos i_c/v_1 = \frac{2h}{v_1}\bigl(\sin^2 i_c/\cos
i_c + \cos i_c\bigr) = 2h/(v_1\cos i_c)$, $\sin i_c = v_1/v_2$ का उपयोग करके: समान। अतिपरवलय की ढलान
$t_r'(x) = x/(v_1\sqrt{x^2 + 4h^2})$ है; $x_*$ पर वह $\sin i_c/v_1 = 1/v_2$ है, जो हेड-तरंग रेखा की ढलान है। चूँकि
$t_r'' > 0$ है, अतिपरवलय उत्तल है और अपनी हर स्पर्श-रेखा के ऊपर रहता है: $x > x_*$ के लिए $t_r > t_{\text{हेड}}$। क्रांतिक
दूरी पर परावर्तन सीमा से ठीक $i_c$ पर टकराता है और वही \omterm{def:b1:layered-earth:head-wave}{हेड तरंग} का जन्म है; उसके आगे \omterm{def:b1:layered-earth:head-wave}{हेड तरंग}, जो तेज़ माध्यम में चलने से
लाभ पाती है, पहले पहुँचती है।
\end{solution}

\begin{solution}{exo:b1:layered-earth:11}
$m(r) = \tfrac43\pi\rho r^3$, अतः $g = Gm/r^2 = \tfrac43\pi G\rho r$। तब
$P(r) = \int_r^R \rho\,\tfrac43\pi G\rho s\,\mathrm ds = \tfrac23\pi G\rho^2(R^2
- r^2)$। केंद्र पर $P = \tfrac23\pi\times\num{6.674e-11}\times5513^2\times
(\num{6.371e6})^2 = \qty{172}{GPa}$, PREM के \qty{363.85}{GPa} का \qty{47}{\percent}: द्रव्यमान को केंद्र की ओर
संकेंद्रित करने से केंद्रीय दाब लगभग दोगुना हो जाता है।
\end{solution}

\begin{solution}{exo:b1:layered-earth:12}
जहाँ $z$ गहराई है, $\mathrm d\rho/\mathrm dz = \rho g/\Phi$; यदि $g/\Phi$ एक स्थिरांक $a$ है, तो $\mathrm d(\ln\rho)/\mathrm dz = a$
और $\rho = \rho_0
e^{a(z - z_0)}$। यहाँ $a(z - z_0) = 7.9\times\num{2258.5e3}/\num{87e6} =
0.2051$ और $\rho = 9.90\times e^{0.2051} = 9.90\times1.2276 =
\qty{12.15}{g.cm^{-3}}$, PREM के $12.17$ से \qty{0.2}{\percent} नीचे: \omterm{def:b1:layered-earth:layers}{बाह्य कोर} में घनत्व की पूरी वृद्धि स्व-संपीड़न है।
\end{solution}

\begin{solution}{pb:b1:layered-earth:1}
\textbf{1.} $V = \tfrac43\pi(\num{6.371e6})^3 = \qty{1.0832e21}{m^3}$.

\textbf{2.} $\tfrac43\pi(r_2^3 - r_1^3)$: \omterm{def:b1:layered-earth:layers}{भूपर्पटी}, ऊपरी \omterm{def:b1:layered-earth:layers}{प्रावार} और संक्रमण-क्षेत्र
\qty{3.0706e20}{m^3}; निचला \omterm{def:b1:layered-earth:layers}{प्रावार} \qty{5.9961e20}{m^3}; \omterm{def:b1:layered-earth:layers}{बाह्य कोर}
\qty{1.6890e20}{m^3}; \omterm{def:b1:layered-earth:layers}{आंतरिक कोर} \qty{7.634e18}{m^3}।

\textbf{3.} \qty{28.35}{\percent}, \qty{55.36}{\percent}, \qty{15.59}{\percent}
और \qty{0.70}{\percent}: निचला \omterm{def:b1:layered-earth:layers}{प्रावार} सबसे बड़ा कोश है, पृथ्वी के आयतन के आधे से अधिक।

\textbf{4.} $15.59 + 0.70 = \qty{16.3}{\percent}$.

\textbf{5.} \omterm{def:b1:layered-earth:layers}{भूपर्पटी}, ऊपरी \omterm{def:b1:layered-earth:layers}{प्रावार} और संक्रमण-क्षेत्र $\num{3.0706e20}\times3570 =
\qty{1.0962e24}{kg}$; निचला \omterm{def:b1:layered-earth:layers}{प्रावार} \qty{2.9381e24}{kg}; \omterm{def:b1:layered-earth:layers}{बाह्य कोर}
\qty{1.8410e24}{kg}; \omterm{def:b1:layered-earth:layers}{आंतरिक कोर} \qty{9.841e22}{kg}।

\textbf{6.} \qty{5.9737e24}{kg}, मापे गए \qty{5.972e24}{kg} से \qty{0.03}{\percent} अधिक।

\textbf{7.} $\num{5.9737e24}/\num{1.0832e21} = \qty{5515}{kg.m^{-3}}$.

\textbf{8.} $\num{9.841e22}/\num{5.9737e24} = \qty{1.65}{\percent}$.

\textbf{9.} $(1.0962 + 2.9381)/5.9737 = \qty{67.5}{\percent}$.

\textbf{10.} कोर आयतन का \qty{16.3}{\percent} और द्रव्यमान का \qty{32.5}{\percent} है, दुगुना: उसका औसत घनत्व,
\qty{10986}{kg.m^{-3}}, \omterm{def:b1:layered-earth:layers}{प्रावार} के लगभग $2.5$ गुना है।

\textbf{11.} एकसमान घनत्व वाला एक कोश देता है
$I = \tfrac{8\pi}{3}\rho\int_{r_1}^{r_2} r^4\,\mathrm dr = \tfrac{8\pi}{15}\rho
(r_2^5 - r_1^5)$।

\textbf{12.} \omterm{def:b1:layered-earth:layers}{भूपर्पटी}, ऊपरी \omterm{def:b1:layered-earth:layers}{प्रावार} और संक्रमण-क्षेत्र \qty{2.676e37}{kg.m^2}; निचला \omterm{def:b1:layered-earth:layers}{प्रावार}
\qty{4.525e37}{kg.m^2}; \omterm{def:b1:layered-earth:layers}{बाह्य कोर} \qty{9.27e36}{kg.m^2}; \omterm{def:b1:layered-earth:layers}{आंतरिक कोर}
\qty{5.87e34}{kg.m^2}।

\textbf{13.} निचला \omterm{def:b1:layered-earth:layers}{प्रावार}: उसके पास लगभग आधा द्रव्यमान है, बड़ी त्रिज्याओं पर। \omterm{def:b1:layered-earth:layers}{आंतरिक कोर} द्रव्यमान का \qty{1.6}{\percent}
रखता है, और वह सब अक्ष के निकट: उसका भाग \qty{0.07}{\percent} है।

\textbf{14.} $I = \qty{8.1345e37}{kg.m^2}$ और $I/MR^2 = \num{8.1345e37}/
(\num{5.9737e24}\times(\num{6.371e6})^2) = 0.3355$।

\textbf{15.} $0.3308$ से \qty{1.4}{\percent} ऊपर, और $0.4$ से काफ़ी नीचे: मॉडल में द्रव्यमान का केंद्र की ओर संकेंद्रण लगभग सही है।

\textbf{16.} हाँ। हर कोश के भीतर वास्तविक घनत्व तल पर शीर्ष की तुलना में अधिक है; समान द्रव्यमान वाला एकसमान कोश उस
द्रव्यमान का बहुत अधिक भाग अपनी बाहरी त्रिज्या के निकट रखता है, जहाँ वह $r^2$ से गिना जाता है, और इसलिए $I$ को अधिक आँकता है।

\textbf{17.} $g = GM/R^2 = \num{6.674e-11}\times\num{5.9737e24}/
(\num{6.371e6})^2 = \qty{9.822}{m.s^{-2}}$.

\textbf{18.} कोर का द्रव्यमान $\num{1.8410e24} + \num{9.841e22} =
\qty{1.9394e24}{kg}$ है, और $g = G m/r_c^2 = \qty{10.69}{m.s^{-2}}$।

\textbf{19.} $(\num{1.0962e24} + \num{2.9381e24})/(\num{3.0706e20} +
\num{5.9961e20}) = \num{4.0343e24}/\num{9.0667e20} = \qty{4450}{kg.m^{-3}}$.

\textbf{20.} $P \approx 4450\times10\times\num{2.891e6} = \qty{128.6}{GPa}$,
PREM के \qty{135.75}{GPa} से \qty{5}{\percent} नीचे।

\textbf{21.} $\rho_c = \num{1.9394e24}/\num{1.7653e20} = \qty{10986}{kg.m^{-3}}$,
और $\Delta P = \tfrac23\pi\times\num{6.674e-11}\times10986^2\times
(\num{3.48e6})^2 = \qty{204.3}{GPa}$।

\textbf{22.} $128.6 + 204.3 = \qty{332.9}{GPa}$, PREM के \qty{363.85}{GPa} से \qty{8.5}{\percent} नीचे। दोनों अनुमान हर क्षेत्र
के भीतर गहराई के साथ घनत्व की वृद्धि को अनदेखा करते हैं; वह अधिक द्रव्यमान को गहराई में रखती है, नीचे उतरते हुए गुरुत्व
को ऊँचा रखती है, और $\rho g$, यानी दाब-प्रवणता, को बढ़ाती है।

\textbf{23.} कुल द्रव्यमान (प्रश्न 6), जड़त्व-आघूर्ण-गुणांक (प्रश्न 15) और पृष्ठ-गुरुत्व (प्रश्न 17, मापे गए औसत
\qty{9.82}{m.s^{-2}} से) पृथ्वी के मापों से जाँचे जाते हैं; कोर पर गुरुत्व और दाबों की तुलना केवल PREM से हो सकती है,
जो स्वयं एक मॉडल है।

\textbf{24.} कोर पृथ्वी के द्रव्यमान का \textbf{\qty[text-series-to-math=true]{32.5}{\percent}} रखता है।
\end{solution}
